\caption{TreeLoom. \textbf{a:} Three hatched intervals of an illustrative simulated ARG yield the matched trees in \textbf{b}. Depth is layout only, with all tips at time zero. Faint bridges join the same recorded ARG node. The magnified split $R$ represents one biological recombination stored as two paired simulator records, with one child lineage and two older parents. Unary nodes are collapsed for display. Variant-site marks correspond to the allele-strip columns. Every tip has observed alleles. Four references supply ancestry badges and one-hot label vectors, while query ancestry is unknown. The attached ancestor vector contains age $\tau$ and zero sampled/root flags. Ancestry colors occur only on labels and outputs. \textbf{c:} Two complete Up--Down GNN traversals use the same weights, indicated by the dashed bracket. The solid upper arrow adds the original node features. Node updates use changed descendant groups and each node's own state. Sample memory preserves every sampled row. Its bidirectional scans decay with physical span $\Delta x$ through $\rho=\exp(-\Delta x/\ell)$ and accumulate current context at $+$. Directional histories are averaged and pooled within the current cut group ($\mu$), centered by current context and projected into that group's representative before traversal two. Matrices retain all $S$ sampled rows, with $q_1$ highlighted and its local vector $e_{tq}$ exposed. \textbf{d:} A separate one-layer BiGRU processes each aligned haplotype along the genome. The chromosome glyph supplies interval geometry. $\Vert$ concatenates local and recurrent context before the classifier. $N_t$ counts native nodes and $d$ is embedding width. Feature values and predictions are schematic. The downward arrow gives the interval-wise argmax, illustrating stable WHG ancestry despite changing surrounding genealogies.}
